\section{Ramification and the complete parameter boundary}
\label{sec:ramified-boundary-v170}
Let $a\geq2$ and retain the marked family and its Hilbert graph $T_a$
from Theorem~\ref{thm:chain-model-v169}. We work over an algebraically
closed field $k$ of characteristic zero. For positive integers
$\rho,\sigma$ consider the finite flat coefficient map
\[
 S'=\Spec k[u,v]\longrightarrow S_a,\qquad b=u^\rho,\quad c=v^\sigma.
\]
Its pullback family is $[t^a:u^\rho s^a:v^\sigma st^{a-1}]$. Put
\[
 Y=T_a\times_{S_a}S',\qquad Z=Y^\nu,\qquad
 \nu:Z\longrightarrow Y,\quad \pi:Z\longrightarrow S'.
\]
Here $Y$ is the retained-coefficient Hilbert graph, whereas $Z$ is its
normalization. The closed fibre $F=\pi^{-1}(0)$ is a scheme of
\emph{parameters}. It is not one of the embedded curves parameterized
by the Hilbert graph. This distinction matters: the conductor of
$\nu$, the nilradical of $F$, and the punctual nilradical of a universal
curve fibre are three different objects.

For $1\leq i\leq a$ define
\begin{equation}\label{eq:ramification-data-v170}
 g_i=\gcd(\rho,i\sigma),\quad
 n_i=(A_i,B_i)=\left(\frac{i\sigma}{g_i},\frac{\rho}{g_i}\right),
 \quad e_i=\frac{\rho\sigma}{g_i},\quad
 \ell_i=\min(A_i,B_i).
\end{equation}
Also set $n_0=(0,1)$ and $n_{a+1}=(1,0)$, ordered clockwise, and
$\Delta_i=|\det(n_i,n_{i+1})|$ for $0\leq i\leq a$.

\begin{theorem}[Ramified Hilbert boundary]\label{thm:ramified-boundary-v170}
The following assertions hold for every $a,\rho,\sigma$ as above.
\begin{enumerate}
\item The scheme $Y$ is integral, is a local complete intersection,
and is the full schematic Hilbert graph of the pulled-back family.
Its normalization $Z$ is the normal toric surface with maximal cones
$\langle n_i,n_{i+1}\rangle$, $0\leq i\leq a$. In particular its
exceptional prime divisors $D_1,\ldots,D_a$ are copies of $\Pj^1$.
\item The fibre $F$ is Cohen--Macaulay of pure dimension one, with
no embedded associated points. Its reduction is a nodal chain of
$a$ projective lines, and its generic length along $D_i$ is $\ell_i$.
Its nilradical $\mathcal N_F$ satisfies
\[
 \mathcal N_F^{L}=0,\qquad L=\max_i\ell_i;
 \quad \mathcal N_F^{L-1}\ne0\quad\text{if }L>1.
\]
It is reduced exactly when, for every $i$, either $\rho\mid i\sigma$
or $i\sigma\mid\rho$.
\item On the open orbit of $D_i$, the normalization maps to the
$i$th exceptional curve of $Y$ with residue degree $g_i$. The finite
map $Z\to T_a$ has ramification index $e_i$ there. A general closed
point of that exceptional curve of $Y$ has exactly $g_i$ distinct
normalization lifts. All torus-fixed boundary points have a single
normalization lift.
\item The conductor, regarded as an ideal on $Z$, is the divisorial
ideal
\begin{equation}\label{eq:conductor-divisor-v170}
 \mathfrak c_{Z/Y}=\OO_Z\left(-\sum_{i=1}^a\kappa_iD_i\right),
 \qquad
 \kappa_i=\frac{\rho i\sigma-\rho-i\sigma}{g_i}+1.
\end{equation}
The notation denotes a reflexive ideal and does not assert that the
divisor is Cartier. There is no additional condition supported only
at a boundary crossing. In particular, $Y$ is normal if and only if
$\rho=1$.
\end{enumerate}
The theorem concerns the complete boundary of this ramified marked
family, not the complete fibre of the full coefficient graph over
$B_a$. None of its multiplicities depends on a choice of Hilbert
embedding degree.
\end{theorem}

\subsection{The normalized Rees algebra}
For $m\geq m_a$ let $E_j$ and $C$ be the integers of
Theorem~\ref{thm:chain-model-v169}, and put
\[
 I=\prod_{j=1}^a(u^\rho,v^{j\sigma})^{E_j},\qquad
 \eta_i=\sum_{j=1}^a E_j\min(i,j).
\]
The raw evaluation ideal on $S'$ is $u^{\rho C}I$.

\begin{proposition}[All powers and all normalization charts]
\label{prop:ramified-rees-v170}
The integral closure of the Rees algebra of $I$ in $k(u,v)(Q)$ is
\begin{equation}\label{eq:ramified-rees-v170}
 \overline{\mathcal R(I)}
 =k\left[u^\alpha v^\beta Q^n:
 \begin{array}{l}
 \alpha,\beta,n\in\mathbb Z_{\geq0},\\
 i\sigma\alpha+\rho\beta\geq n\rho\sigma\eta_i
                  \quad(1\leq i\leq a)
 \end{array}\right].
\end{equation}
Thus its degree-$n$ part is $\overline{I^n}Q^n$, including $n=0$.
For the raw ideal one replaces $\alpha\geq0$ by
$\alpha\geq n\rho C$ and the right side of the $i$th inequality by
$n\rho\sigma(\eta_i+iC)$. The affine chart for the cone
$\tau_i=\langle n_i,n_{i+1}\rangle$ is exactly
\begin{equation}\label{eq:normal-chart-v170}
 Z_i=\Spec k[\tau_i^\vee\cap\mathbb Z^2].
\end{equation}
Adjacent charts glue by inverting the characters of their shared
face in the common Laurent ring $k[u^{\pm1},v^{\pm1}]$.
\end{proposition}
\begin{proof}
The coefficient map is free with basis $u^jv^k$, $0\leq j<\rho$,
$0\leq k<\sigma$. Flat base change of the Rees algebra consequently
gives $Y=\operatorname{Bl}_I S'$; see the classical flat-base-change
lemma for blow-ups \cite{Stacks169}. The open $uv\ne0$ is
schematically dense by flat base change from $T_a$. Its coordinate
ring is a domain, so $Y$ is integral. The same density identifies
$Y$ with the closure of the pulled-back Hilbert section. In particular
no separate vertical component has been discarded in this equality.

The Newton polygon of $I$ is the image, under
$(x,y)\mapsto(\rho x,\sigma y)$, of the polygon calculated in
Proposition~\ref{prop:rees-v169}. Its compact supporting inequalities
are precisely those in \eqref{eq:ramified-rees-v170}. All compact
edges have positive length. The semigroup of the Rees algebra
contains the two degree-zero basis vectors and a degree-one vector,
so its group is the whole $\mathbb Z^3$. Saturating that semigroup
therefore gives the integral lattice points in its rational cone.
Indeed, a rational positive combination becomes an element of the
original semigroup after clearing denominators; its monomial is
integral. Conversely every supporting valuation is nonnegative on
an integral element. The saturated semigroup ring is normal. This
is the standard semigroup-normalization theorem
\cite[Proposition~5 and Section~2.6]{GonzalezTeissier170}, here applied
to the explicit polygon. It proves the formula in every degree.

The normal fan is unchanged by the positive edge lengths $E_j$.
Its primitive compact rays are exactly $n_i$, giving
\eqref{eq:normal-chart-v170}. Normalization commutes with localization,
so the stated overlaps give the global normalization, not unrelated
normalizations of completed rings. The shift by $u^{n\rho C}$ proves
the raw formula and does not change Proj.
\end{proof}

\begin{proposition}[Singularities and divisor data]
\label{prop:ramified-singularities-v170}
The singularities of $Z$ are precisely those torus-fixed points whose
corresponding determinant in the following list is greater than one:
\begin{equation}\label{eq:cone-indices-v170}
 \Delta_0=\frac{\sigma}{g_1},\qquad
 \Delta_i=\frac{\rho\sigma}{g_i g_{i+1}}\ (1\leq i<a),\qquad
 \Delta_a=\frac{\rho}{g_a}.
\end{equation}
The affine chart has cyclic divisor class group
$\mathbb Z/\Delta_i\mathbb Z$. Its cyclic quotient type is determined
without ambiguity by the two primitive rays in
\eqref{eq:normal-chart-v170}: orient them positively, extend the
first to a lattice basis, and write the second as $(p,\Delta_i)$
in that basis. The chart is $\A^2/\mu_{\Delta_i}$ with weights
$(-p,1)$, up to swapping the coordinates. For $\Delta_i=1$ it is smooth.
Moreover
\[
 \operatorname{div}(u)=U'+\sum_i A_iD_i,
 \qquad \operatorname{div}(v)=V'+\sum_i B_iD_i,
 \qquad K_{Z/S'}=\sum_i(A_i+B_i-1)D_i,
\]
where $U',V'$ are the strict transforms of the axes. Rational
intersection numbers are
\[
 D_iD_{i+1}=\Delta_i^{-1},\qquad
 D_i^2=-\frac{|\det(n_{i-1},n_{i+1})|}{\Delta_{i-1}\Delta_i},
 \qquad D_iD_j=0\quad(|i-j|>1).
\]
\end{proposition}
\begin{proof}
The determinant formulas follow by substitution. For a cone generated
by primitive vectors $r_1,r_2$, its ray sublattice has index
$|\det(r_1,r_2)|$. The corresponding toric chart in that sublattice
is $\A^2$. The lattice quotient is cyclic because the first ray is
primitive; in the indicated basis its generator acts with weights
$(-p,1)$. Primitivity of the second ray implies
$\gcd(p,\Delta_i)=1$, so there are no pseudoreflections and the
quotient is smooth exactly for index one. The divisor class group
is the cokernel of the map from the character lattice to the free
group on the two rays; its Smith form is $(1,\Delta_i)$.
Normal surface charts are regular away from their closed toric point.

A character has order its pairing with a primitive ray. This proves
the two divisor formulas. The torus form
$(du/u)\wedge(dv/v)$ has simple poles on each invariant divisor,
which proves the relative canonical formula. The two adjacent
invariant divisors meet with rational intersection $1/\Delta_i$:
pull back to the finite affine cover, where the axes meet with
intersection one and the covering degree is $\Delta_i$. Intersecting
the principal-character relations with $D_i$ then gives the displayed
self-intersection. These calculations also specify a toric resolution
by subdivision into unimodular cones; its curves and intersection
numbers are obtained by the same determinant rule.
\end{proof}

\subsection{The scheme fibre, not only its cycle}
The following elementary lattice observation rules out hidden
zero-dimensional embedded components in the parameter fibre.

\begin{lemma}[The two coordinate ideal on a subdivided surface]
\label{lem:coordinate-ideal-v170}
Let $\tau$ be a two-dimensional cone inside the first quadrant with
primitive boundary vectors $l,h$, and suppose $\tau$ is not the
whole first quadrant. In $A=k[\tau^\vee\cap\mathbb Z^2]$ one has
\begin{equation}\label{eq:fibre-divisorial-v170}
 (u,v)=\left\langle u^xv^y:
 \langle l,(x,y)\rangle\geq\min(l_1,l_2),\quad
 \langle h,(x,y)\rangle\geq\min(h_1,h_2)\right\rangle_k.
\end{equation}
The exponents may be negative; the inequalities imply membership in
$A$. The ideal is divisorial, and $A/(u,v)$ is Cohen--Macaulay of
dimension one whenever its support is nonempty.
\end{lemma}
\begin{proof}
If both rays are on the same side of $(1,1)$, either $u/v$ or $v/u$
is regular on the chart, and the ideal is principal. Otherwise write
$l=(l_1,l_2)$, $h=(h_1,h_2)$ with
$l_1<l_2$ and $h_1>h_2$; equality cases belong to the preceding
argument. For a lattice point satisfying the displayed inequalities,
put $L=x+y$. If $L\leq0$ and $l_1,h_2>0$, the first inequality
forces $y>0$, while the second forces $y<0$. If one of $l_1,h_2$
is zero, its inequality instead forces the weak sign, and the other
strict inequality gives the contradiction. They cannot both vanish,
since the original quadrant was excluded. Thus $L\geq1$.

If neither $(x-1,y)$ nor $(x,y-1)$ belonged to the dual cone, the
only possible failed inequalities would be
\[
 l_1L+(l_2-l_1)y<l_2,
 \qquad h_1L-(h_1-h_2)y<h_1.
\]
The first implies $y\leq0$ and the second $y>0$, since $L\geq1$
and $y$ is integral. This is impossible. Thus the monomial is a
multiple of $u$ or $v$. The reverse inclusion is immediate.

The formula is the intersection of the symbolic valuation ideals
of the invariant height-one primes. All noninvariant height-one
conditions are supplied by membership in the normal ring $A$.
Consequently this rank-one ideal is reflexive. A reflexive module
over a normal surface has depth two at closed points. The exact
sequence with middle term $A$ now gives depth at least one for its
quotient, as asserted.
\end{proof}

\begin{proof}[Proof of the fibre assertions in Theorem~\ref{thm:ramified-boundary-v170}]
Every maximal cone of $Z$ is a proper subcone of the quadrant, so
Lemma~\ref{lem:coordinate-ideal-v170} applies. Globally it says
\begin{equation}\label{eq:complete-fibre-ideal-v170}
 (u,v)\OO_Z=\OO_Z\left(-\sum_i\ell_iD_i\right).
\end{equation}
The fibre thus has no embedded associated points. Its generic length
is the indicated DVR order $\ell_i$. Reducing the ideal replaces
each positive order by one, so the reduction is the chain of
invariant projective lines. At a cone with two exceptional edges,
the reduced union has ring $k[z_1,z_2]/(z_1z_2)$: on the cyclic
quotient cover the two axes have invariant coordinate rings generated
by the respective $\Delta_i$th powers. Thus the crossings really
are nodes, even when the ambient surface is singular.

The radical ideal has order one at every $D_i$. Its $L$th power lies
in \eqref{eq:complete-fibre-ideal-v170} by the divisorial description,
whereas localization at a divisor of length $L$ shows nonvanishing
of the preceding nilradical power if $L>1$. Finally
$\min(\rho,i\sigma)/g_i=1$ is equivalent to one of the two integers
$\rho,i\sigma$ dividing the other. This proves all claims, including
the exact nilpotence order rather than only a bound from a cycle.
\end{proof}
